% ============================================================
\chapter{Paradigme Bay\'esien}
\label{ch:paradigme}
% ============================================================

En 1763, deux ans apr\`es sa mort, le r\'ev\'erend Thomas Bayes voit son
essai publi\'e \`a titre posthume par la Royal Society de Londres. L'article
propose une m\'ethode pour inverser les probabilit\'es~: connaissant les
donn\'ees observ\'ees, que peut-on dire de la cause qui les a produites~?
Cette question, d'apparence innocente, va engendrer l'un des d\'ebats les
plus durables de toute la statistique. D'un c\^ot\'e, les
\emph{fr\'equentistes}, pour qui les param\`etres sont des constantes fixes
et inconnues. De l'autre, les \emph{bay\'esiens}, pour qui les param\`etres
sont des variables al\'eatoires dot\'ees d'une distribution refl\'etant notre
\'etat de connaissance. Ce chapitre pr\'esente le paradigme bay\'esien~: sa
logique, ses forces, et les raisons pour lesquelles il conna\^it
aujourd'hui un renouveau spectaculaire.

\begin{intuition}[title=\textbf{Id\'ee centrale}]
Le paradigme bay\'esien traite le param\`etre $\theta$ comme une variable
al\'eatoire et encode toute l'incertitude dans des distributions de
probabilit\'e. L'inf\'erence consiste \`a passer de la loi a priori
$\pi(\theta)$ \`a la loi a posteriori $\pi(\theta \mid x)$ via le
th\'eor\`eme de Bayes.
\end{intuition}

% ------------------------------------------------------------
\section{Introduction et motivation}
% ------------------------------------------------------------

L'approche fr\'equentiste consid\`ere le param\`etre $\theta$ comme une
constante fixe inconnue et \'evalue les proc\'edures d'inf\'erence par
leurs propri\'et\'es r\'ep\'et\'ees sur des \'echantillons hypoth\'etiques.
L'approche bay\'esienne, en revanche, mod\'elise $\theta$ comme une
variable al\'eatoire dont la distribution refl\`ete notre \'etat de
connaissance.

Cette diff\'erence de philosophie entra\^ine des cons\'equences profondes~:
\begin{itemize}
  \item L'incertitude est directement quantifi\'ee par la loi a posteriori.
  \item L'information a priori est formellement int\'egr\'ee.
  \item Toute d\'ecision optimale peut \^etre d\'eduite du posterior.
\end{itemize}

% ------------------------------------------------------------
\section{Fondements axiomatiques}
% ------------------------------------------------------------

\subsection{Probabilit\'e subjective}

\begin{definition}[Probabilit\'e subjective]
Une \textbf{probabilit\'e subjective} est une fonction
$\PP : \mathcal{F} \to [0,1]$ satisfaisant les axiomes de Kolmogorov,
interpr\'et\'ee comme le degr\'e de croyance d'un agent rationnel dans
la r\'ealisation d'un \'ev\'enement.
\end{definition}

Les axiomes de coh\'erence de de Finetti montrent qu'un agent dont les
paris sont coh\'erents (pas de \textit{Dutch book}) se comporte
n\'ecessairement comme s'il utilisait des probabilit\'es satisfaisant
les axiomes de Kolmogorov.

\begin{theoreme}[Th\'eor\`eme de repr\'esentation de de Finetti]
Si $(X_1, X_2, \ldots)$ est une suite \'echangeable de variables
al\'eatoires \`a valeurs dans $\{0,1\}$, alors il existe une mesure de
probabilit\'e $\mu$ sur $[0,1]$ telle que~:
\[
  \PP(X_1 = x_1, \ldots, X_n = x_n)
  = \int_0^1 \theta^{s_n}(1-\theta)^{n - s_n}\,d\mu(\theta),
\]
o\`u $s_n = \sum_{i=1}^n x_i$.
\end{theoreme}

\begin{remarque}
Ce th\'eor\`eme justifie le mod\`ele bay\'esien~: l'\'echangeabilit\'e
(hypoth\`ese plus faible que l'ind\'ependance) implique l'existence
d'un param\`etre latent $\theta$ avec une distribution a priori $\mu$.
\end{remarque}

\subsection{Axiomes de Savage}

\begin{theoreme}[Th\'eor\`eme de Savage, 1954]
Sous les sept axiomes de Savage (ordonnancement, \textit{sure-thing
principle}, etc.), il existe~:
\begin{enumerate}
  \item une mesure de probabilit\'e subjective $\PP$ sur l'espace des
        \'etats,
  \item une fonction d'utilit\'e $u$ unique \`a transformation affine
        pr\`es,
\end{enumerate}
telles que l'agent pr\'ef\`ere l'action $a$ \`a l'action $b$ si et
seulement si $\E_\PP[u(a)] > \E_\PP[u(b)]$.
\end{theoreme}

% ------------------------------------------------------------
\section{Le th\'eor\`eme de Bayes}
% ------------------------------------------------------------

\begin{theoreme}[Th\'eor\`eme de Bayes]
\label{thm:bayes}
Soit $\theta \in \Theta$ un param\`etre de loi a priori $\pi(\theta)$
et $x = (x_1, \ldots, x_n)$ un \'echantillon de densit\'e
$f(x \mid \theta)$. La loi a posteriori est~:
\[
  \pi(\theta \mid x)
  = \frac{f(x \mid \theta)\,\pi(\theta)}{m(x)},
  \qquad
  m(x) = \int_\Theta f(x \mid \theta)\,\pi(\theta)\,d\theta.
\]
\end{theoreme}

\begin{proof}
Par la r\`egle de Bayes sur les densit\'es conditionnelles~:
\[
  \pi(\theta \mid x)
  = \frac{f(x, \theta)}{f(x)}
  = \frac{f(x \mid \theta)\,\pi(\theta)}
         {\int_\Theta f(x \mid \theta)\,\pi(\theta)\,d\theta}.
\]
\end{proof}

\begin{formules}[title=\textbf{Composantes du th\'eor\`eme de Bayes}]
\begin{align*}
  \text{Prior} &: \quad \pi(\theta) \\
  \text{Vraisemblance} &: \quad L(\theta \mid x) = f(x \mid \theta) \\
  \text{Evidence} &: \quad m(x) = \int_\Theta L(\theta \mid x)\,\pi(\theta)\,d\theta \\
  \text{Posterior} &: \quad \pi(\theta \mid x) \propto L(\theta \mid x)\,\pi(\theta)
\end{align*}
\end{formules}

% ------------------------------------------------------------
\section{Choix de la loi a priori}
% ------------------------------------------------------------

\subsection{Types de priors}

\begin{definition}[Prior informatif]
Un \textbf{prior informatif} est une loi $\pi(\theta)$ qui concentre
sa masse sur une r\'egion restreinte de $\Theta$, refl\'etant une
connaissance pr\'ealable substantielle.
\end{definition}

\begin{definition}[Prior vague (non informatif)]
Un \textbf{prior vague} ou \textbf{non informatif} est une loi
$\pi(\theta)$ qui cherche \`a laisser les donn\'ees dominer
l'inf\'erence. Exemples~: prior uniforme, prior de Jeffreys.
\end{definition}

\begin{definition}[Prior de Jeffreys]
Le \textbf{prior de Jeffreys} est d\'efini par~:
\[
  \pi_J(\theta) \propto \sqrt{\det I(\theta)},
\]
o\`u $I(\theta)$ est la matrice d'information de Fisher.
\end{definition}

\begin{proposition}[Invariance du prior de Jeffreys]
Le prior de Jeffreys est invariant par reparam\'etrisation~: si
$\phi = g(\theta)$ est une bijection diff\'erentiable, alors
$\pi_J(\phi) \propto \sqrt{\det I_\phi(\phi)}$.
\end{proposition}

\begin{proof}
Soit $\phi = g(\theta)$. Par la r\`egle de changement de variable et
la transformation de l'information de Fisher~:
\[
  I_\phi(\phi) = \left(\frac{d\theta}{d\phi}\right)^2 I(\theta),
\]
d'o\`u $\sqrt{I_\phi(\phi)} = \abs{\frac{d\theta}{d\phi}} \sqrt{I(\theta)}$,
ce qui correspond exactement au jacobien du changement de variable.
\end{proof}

\begin{exemple}
Pour $X_1, \ldots, X_n \overset{\text{iid}}{\sim} \mathcal{N}(\mu, \sigma^2)$
avec $\sigma^2$ connu, l'information de Fisher pour $\mu$ est
$I(\mu) = n/\sigma^2$, constante. Le prior de Jeffreys est donc
$\pi_J(\mu) \propto 1$, c'est-\`a-dire le prior uniforme (impropre)
sur $\R$.
\end{exemple}

\subsection{Prior de r\'ef\'erence}

\begin{definition}[Prior de r\'ef\'erence de Bernardo]
Le \textbf{prior de r\'ef\'erence} maximise l'information attendue de
Kullback--Leibler entre la loi a priori et la loi a posteriori~:
\[
  \pi^*(\theta) = \argmax_{\pi} \E_x\left[\KL\big(\pi(\theta \mid x) \| \pi(\theta)\big)\right].
\]
\end{definition}

% ------------------------------------------------------------
\section{Exemples fondamentaux}
% ------------------------------------------------------------

\begin{exemple}
\textbf{Mod\`ele Bernoulli--B\^eta.}
Soit $X_1, \ldots, X_n \overset{\text{iid}}{\sim} \text{Ber}(\theta)$
et $\theta \sim \text{Be}(a, b)$. Posons $s = \sum x_i$. Alors~:
\begin{align*}
  \pi(\theta \mid x)
  &\propto \theta^s (1-\theta)^{n-s} \cdot \theta^{a-1}(1-\theta)^{b-1} \\
  &= \theta^{a+s-1}(1-\theta)^{b+n-s-1}.
\end{align*}
Donc $\theta \mid x \sim \text{Be}(a + s, \, b + n - s)$.
\end{exemple}

\begin{exemple}
\textbf{Mod\`ele Normal--Normal (variance connue).}
Soit $X_1, \ldots, X_n \overset{\text{iid}}{\sim} \mathcal{N}(\mu, \sigma^2)$
avec $\sigma^2$ connu, et $\mu \sim \mathcal{N}(\mu_0, \tau_0^2)$.
La loi a posteriori est~:
\[
  \mu \mid x \sim \mathcal{N}(\mu_n, \tau_n^2),
\]
o\`u
\[
  \tau_n^2 = \left(\frac{1}{\tau_0^2} + \frac{n}{\sigma^2}\right)^{-1}, \qquad
  \mu_n = \tau_n^2 \left(\frac{\mu_0}{\tau_0^2} + \frac{n\bar{x}}{\sigma^2}\right).
\]
\end{exemple}

\begin{formules}[title=\textbf{Moyenne a posteriori comme moyenne pond\'er\'ee}]
\[
  \mu_n = w \cdot \mu_0 + (1-w) \cdot \bar{x},
  \qquad w = \frac{\sigma^2/n}{\tau_0^2 + \sigma^2/n}.
\]
Quand $n \to \infty$, $w \to 0$ et $\mu_n \to \bar{x}$~: le posterior
converge vers l'estimateur du maximum de vraisemblance.
\end{formules}

% ------------------------------------------------------------
\section{Propri\'et\'es asymptotiques}
% ------------------------------------------------------------

\begin{theoreme}[Th\'eor\`eme de Bernstein--von Mises]
Sous des conditions de r\'egularit\'e, la loi a posteriori converge
en loi vers une gaussienne centr\'ee sur le maximum de vraisemblance~:
\[
  \pi(\theta \mid x_1, \ldots, x_n)
  \xrightarrow{d}
  \mathcal{N}\!\left(\hat{\theta}_{\text{MLE}},\;
  \frac{1}{n}\,I(\theta_0)^{-1}\right),
\]
o\`u $\theta_0$ est la vraie valeur du param\`etre.
\end{theoreme}

\begin{remarque}
Ce r\'esultat montre que, pour de grands \'echantillons, le choix du
prior a un impact n\'egligeable. Autrement dit, les approches
bay\'esienne et fr\'equentiste convergent asymptotiquement.
\end{remarque}

\begin{theoreme}[Consistance a posteriori]
Soit $\theta_0$ la vraie valeur. Sous des conditions de r\'egularit\'e
(identifiabilit\'e, support du prior contenant $\theta_0$),
pour tout voisinage $U$ de $\theta_0$~:
\[
  \pi(\theta \in U \mid x_1, \ldots, x_n) \xrightarrow{p} 1
  \quad \text{quand } n \to \infty.
\]
\end{theoreme}

% ------------------------------------------------------------
\section{Principe de vraisemblance}
% ------------------------------------------------------------

\begin{definition}[Principe de vraisemblance]
Deux exp\'eriences produisant la m\^eme fonction de vraisemblance
$L(\theta \mid x)$ (\'a une constante multiplicative pr\`es) doivent
conduire \`a la m\^eme inf\'erence sur $\theta$.
\end{definition}

\begin{proposition}
L'inf\'erence bay\'esienne satisfait automatiquement le principe de
vraisemblance, car la loi a posteriori ne d\'epend de $x$ qu'\`a
travers $L(\theta \mid x)$.
\end{proposition}

\begin{attention}[title=\textbf{Violation fr\'equentiste}]
Certaines proc\'edures fr\'equentistes (comme les $p$-valeurs) violent
le principe de vraisemblance car elles d\'ependent de l'espace
\'echantillonnal et du plan d'exp\'erience, pas seulement des donn\'ees
observ\'ees.
\end{attention}

% ------------------------------------------------------------
\section{Impl\'ementation Python}
% ------------------------------------------------------------

\begin{codeblock}[title=\textbf{Mod\`ele Bernoulli--B\^eta avec PyMC}]
\begin{minted}{python}
import pymc as pm
import arviz as az
import numpy as np

# Données simulées
np.random.seed(42)
n = 50
theta_true = 0.3
data = np.random.binomial(1, theta_true, size=n)

# Modèle bayésien
with pm.Model() as model_beta_binom:
    # Prior Beta(2, 5)
    theta = pm.Beta("theta", alpha=2, beta=5)
    # Vraisemblance
    y_obs = pm.Bernoulli("y_obs", p=theta, observed=data)
    # Échantillonnage MCMC
    trace = pm.sample(2000, tune=1000, random_seed=42)

# Résumé du posterior
summary = az.summary(trace, var_names=["theta"],
                     hdi_prob=0.95)
print(summary)

# Vérification analytique
a_post = 2 + data.sum()
b_post = 5 + n - data.sum()
print(f"Posterior analytique: Be({a_post}, {b_post})")
print(f"Moyenne analytique: {a_post/(a_post+b_post):.4f}")
\end{minted}
\end{codeblock}

\begin{codeblock}[title=\textbf{Visualisation prior--posterior}]
\begin{minted}{python}
import matplotlib.pyplot as plt
from scipy import stats

fig, ax = plt.subplots(figsize=(8, 5))
theta_grid = np.linspace(0, 1, 500)

# Prior
ax.plot(theta_grid, stats.beta.pdf(theta_grid, 2, 5),
        label="Prior Be(2, 5)", linestyle="--")
# Posterior analytique
ax.plot(theta_grid, stats.beta.pdf(theta_grid, a_post, b_post),
        label=f"Posterior Be({a_post}, {b_post})")
# Vraisemblance (normalisée pour affichage)
s = data.sum()
like = theta_grid**s * (1-theta_grid)**(n-s)
like /= like.max() * 3
ax.plot(theta_grid, like, label="Vraisemblance (norm.)",
        linestyle=":")
ax.axvline(theta_true, color="red", alpha=0.5,
           label=f"Vraie valeur ({theta_true})")
ax.set_xlabel(r"$\theta$")
ax.set_ylabel("Densité")
ax.legend()
ax.set_title("Mise à jour bayésienne: Bernoulli-Bêta")
plt.tight_layout()
plt.savefig("bayes_update.pdf")
\end{minted}
\end{codeblock}

% ------------------------------------------------------------
\section{Exercices}
% ------------------------------------------------------------

\begin{exercice}
Soit $X_1, \ldots, X_n \overset{\text{iid}}{\sim} \text{Poi}(\lambda)$
avec prior $\lambda \sim \text{Ga}(\alpha, \beta)$.
\begin{enumerate}
  \item D\'eterminer la loi a posteriori de $\lambda$.
  \item Calculer la moyenne et la variance a posteriori.
  \item Montrer que la moyenne a posteriori est une moyenne
        pond\'er\'ee de la moyenne a priori et de la moyenne
        empirique.
\end{enumerate}
\end{exercice}

\begin{exercice}
Calculer le prior de Jeffreys pour le mod\`ele $\text{Ber}(\theta)$ et
montrer qu'il s'agit de $\text{Be}(1/2, 1/2)$.
\end{exercice}

\begin{exercice}
D\'emontrer que, dans le mod\`ele Normal--Normal, la pr\'ecision a
posteriori est la somme de la pr\'ecision a priori et de la pr\'ecision
des donn\'ees. Interpr\'eter ce r\'esultat.
\end{exercice}

\begin{exercice}
Consid\'erer deux exp\'eriences~:
\begin{enumerate}
  \item On lance une pi\`ece $n = 12$ fois et on observe $s = 3$ succ\`es.
  \item On lance une pi\`ece jusqu'au troisi\`eme succ\`es et on observe
        $n = 12$ lancers.
\end{enumerate}
Montrer que les vraisemblances sont proportionnelles. En d\'eduire que
l'inf\'erence bay\'esienne est identique, alors que les $p$-valeurs
fr\'equentistes diff\`erent.
\end{exercice}

\begin{exercice}
\textbf{(Num\'erique)} Reprendre le mod\`ele Normal--Normal avec
$\sigma^2 = 1$, $\mu_0 = 0$, $\tau_0^2 = 10$ et un \'echantillon de
taille $n = 5$ simul\'e depuis $\mathcal{N}(2, 1)$.
\begin{enumerate}
  \item Calculer analytiquement le posterior.
  \item V\'erifier avec PyMC.
  \item Tracer le prior, le posterior et la vraisemblance normalis\'ee
        sur un m\^eme graphique.
\end{enumerate}
\end{exercice}
